\textbf{COLUNI 2021 -- Questão 25}\par\medskip Vênus, a joia do céu, era conhecida pelos astrônomos da antiguidade como a Estrela D'álva ou Estrela Vespertina. Vênus, assim chamado por causa da deusa romana do amor e da beleza, é encoberto por grossa camada de nuvens em turbilhões. A temperatura média da superfície de Vênus é cerca de 482 °C. Essa alta temperatura é devida, especialmente, ao efeito estufa, causado pela pesada atmosfera de dióxido de carbono (CO$_2$). Isso torna Vênus mais quente que Mercúrio. Disponível em: http://astro.if.ufrgs.br/solar/venus.htm. Acesso em: 18 nov. 2020. (Adaptado)

\par

Vênus ganhou destaque em 2020, em função da descoberta de atividade microbiana por um grupo internacional de cientistas que publicou os resultados da pesquisa na Revista \textit{Nature Astronomy}.

\par

Observe a tabela abaixo, na qual algumas substâncias encontradas em nosso planeta estão listadas:

\par

\begin{center}
\begin{tabular}{|l|c|c|}
\hline
\textbf{Substância} & \textbf{Ponto de fusão (ºC)} & \textbf{Ponto de ebulição (ºC)} \\
\hline
Fosfina (PH$_3$) & $-132{,}8$ & $-87{,}7$ \\
\hline
Cloreto de sódio (NaCl) & 801 & 1465 \\
\hline
Mercúrio (Hg) & $-38{,}8$ & 356{,}7 \\
\hline
Chumbo (Pb) & 327{,}5 & 1749 \\
\hline
\end{tabular}
\end{center}

\par

Considerando-se a temperatura da superfície do planeta Vênus, os estados físicos da fosfina, do cloreto de sódio, do mercúrio e do chumbo são, respectivamente:\begin{enumerate}
\item sólido, gasoso, líquido, gasoso.
\item gasoso, sólido, gasoso, líquido.
\item gasoso, líquido, gasoso, sólido.
\item líquido, sólido, líquido, líquido.\end{enumerate}
